\subsection{Proceso de ingeniería de requisitos}

\subsubsection{Actividades centrales del marco de Pohl}

Pohl distingue tres actividades centrales: la elicitación, que permite
obtener información de las fuentes, la documentación, que registra
los requisitos y sus decisiones, y la negociación, que busca resolver
diferencias entre los interesados~\cite{pohl2025}.

Durante la elicitación, el equipo puede consultar a quienes realizan
una tarea y revisar los documentos que usan. La documentación
permite conservar qué necesidad se identificó y de dónde proviene.
Si dos actores esperan resultados distintos, la negociación ayuda
a decidir qué se incorpora al alcance y a registrar el motivo de
esa decisión~\cite{pohl2025,swebok2024}.

\textbf{Ejemplo en UteqFit-AI.}
Para revisar RI-01, el equipo puede solicitar a Bienestar
Universitario una descripción del listado de estudiantes habilitados
y consultar al encargado sobre su uso. Después puede documentar
cuándo se entrega el listado y qué información necesita la
importación. Si se pide una actualización distinta a la prevista
por período académico, el equipo tendría que acordar quién la
autoriza y cómo se entrega.

\subsubsection{Actividades transversales}

En el marco de Pohl, la validación comprueba si los requisitos
representan las necesidades de los interesados. La gestión de
requisitos mantiene su identificación, sus relaciones y su historial
de cambios durante el proyecto~\cite{pohl2025,ieee29148}.

La validación permite revisar un borrador con quienes conocen la
operación del gimnasio. Si esa revisión revela un caso omitido,
el equipo puede actualizar el requisito y su criterio de aceptación.
La gestión ayuda a conservar la relación entre el cambio, la fuente
que lo motivó y los requisitos que dependen de él~\cite{pohl2025,swebok2024}.

\textbf{Ejemplo en UteqFit-AI.}
Para revisar la importación del listado, se pueden preparar dos
casos con datos ficticios: un estudiante que figura como habilitado
y otro que no aparece en el archivo. Bienestar Universitario y el
encargado de turno pueden revisar el resultado esperado para cada
caso. Si acuerdan un tratamiento especial para las discrepancias,
el equipo puede registrar la decisión en RI-01 y revisar su
compatibilidad con RR-01.

\subsubsection{Secuencia de trabajo aplicada al PFC}

SWEBOK organiza el trabajo con requisitos en actividades que incluyen
elicitación, análisis, especificación y validación. Estas actividades
pueden repetirse cuando aparece información que obliga a modificar
un requisito~\cite{swebok2024,ieee29148}.

La elicitación permite identificar qué necesita cada actor y qué
información puede aportar. En el análisis, el equipo contrasta esa
información con el alcance del sistema y revisa posibles
contradicciones. En la especificación, convierte las decisiones
en enunciados identificables y verificables. En la validación,
los interesados revisan si los enunciados y los casos de aceptación
reflejan la operación esperada~\cite{swebok2024,ieee29148}.

\textbf{Ejemplo en UteqFit-AI.}
Para el control de ocupación, la secuencia propuesta comienza
con una consulta al encargado sobre el registro de entradas y
salidas. Sigue con el análisis de qué cuenta como visita abierta
y con la especificación de RF-01, RF-02 y RF-06. La revisión
con el responsable del gimnasio puede utilizar casos de entrada,
salida y visita sin cerrar. Como evidencia de aplicación, el
equipo puede adjuntar la nota de la consulta, la versión de los
requisitos revisada y el registro de los resultados cuando se
realice la validación.